\documentclass[11pt,a4paper]{article}
\usepackage{amsmath,amssymb}
\usepackage{listings}
\usepackage{hyperref}
\usepackage[margin=1in]{geometry}
\usepackage{fancyhdr}

\pagestyle{fancy}
\fancyhf{}
\fancyhead[L]{Module 18: Handwriting Recognition}
\fancyhead[R]{C Programming Course}
\fancyfoot[C]{\thepage}

\lstset{
  language=C,
  basicstyle=\ttfamily\small,
  keywordstyle=\bfseries,
  commentstyle=\itshape,
  numbers=left,
  numberstyle=\tiny,
  frame=single,
  breaklines=true,
  tabsize=4
}

\title{Module 18: Handwriting Recognition}
\author{C Programming Course}
\date{}

\begin{document}
\maketitle
\thispagestyle{fancy}

\section{Understanding the MNIST Dataset}
MNIST contains 70,000 grayscale images of handwritten digits (0--9), each
$28 \times 28$ pixels. The dataset is split into 60,000 training and 10,000
test images. Pixel values range from 0 (white) to 255 (black).

\section{Image Processing in C}
Each image is a flat array of 784 bytes. Normalize pixel values to $[0, 1]$:

\begin{lstlisting}
#include <stdio.h>
#include <stdint.h>

#define IMG_SIZE 784

void normalize(const uint8_t *raw, double *out) {
    for (int i = 0; i < IMG_SIZE; i++) {
        out[i] = raw[i] / 255.0;
    }
}
\end{lstlisting}

\section{Building a Digit Recognizer}
A simple fully connected network with one hidden layer:
\begin{itemize}
  \item Input layer: 784 neurons.
  \item Hidden layer: 128 neurons with ReLU activation.
  \item Output layer: 10 neurons with softmax.
\end{itemize}

\begin{lstlisting}
#include <math.h>

void softmax(double *z, int n) {
    double max_val = z[0];
    for (int i = 1; i < n; i++)
        if (z[i] > max_val) max_val = z[i];
    double sum = 0.0;
    for (int i = 0; i < n; i++) {
        z[i] = exp(z[i] - max_val);
        sum += z[i];
    }
    for (int i = 0; i < n; i++)
        z[i] /= sum;
}
\end{lstlisting}

\section{Training and Testing}
\begin{enumerate}
  \item Load and normalize the MNIST data.
  \item Initialize weights with small random values.
  \item Train using mini-batch gradient descent.
  \item Evaluate accuracy on the test set.
\end{enumerate}

Cross-entropy loss for classification:
\[
L = -\sum_{i=0}^{9} y_i \log(\hat{y}_i)
\]

\section{Optimization Techniques}
\begin{itemize}
  \item \textbf{Learning rate decay}: reduce $\eta$ over time.
  \item \textbf{Mini-batch training}: update weights every $k$ samples.
  \item \textbf{Weight initialization}: use He or Xavier initialization.
\end{itemize}

\end{document}
